\scanpage{76}
% archive sheet number: 75
% @ L1–11
par vrillement \uline{\unsure{de deux|des|deux|de}} polygones \uline{d'ordre $n$}, non $2n$ ? \deleted{\unsure{Cett|Cette|Celle}} La première ne serait-elle pas plutôt un revêtement d'ordre deux de la seconde, ramifiée en chacun des pts du centre -- en effet, on voit \deleted{apparemment} \inserted{que} les antipodismes des les $\Delta_D^*$ sont compatibles avec les identifications, -- par passage au quotient, cela va donner sûrement la carte cellulaire envisagée\ldots{}

% @ L11–20
On peut alors, \deleted{qui} après \deleted{\unsure{avoir enlevé|avoir admis}} interpréter la nouvelle carte comme associée à un graphe \inserted{(plongé cellulairement)} dont les sommets sont les pos. stables, qui pour deux pos. $D,D'$ adjacentes, on va maintenant envisager \uline{deux} flèches ou arêtes entre eux, savoir $a_D:D\mathbin{\text{---}}D'$ (de l'\unsure{about|arête} $a$ vers $b'$) et $a_{D'}:D'\mathbin{\text{---}}D$ (de l'\unsure{arête|autr} $b$ vers $a'$), munies chacune d'une orientation -- le balayage de $D$ vers $D'$, et de $D'$ vers $D$, qu'on désidentifie.

% @ R1–8
Maintnt, les arcs de transition (nouveau style) pour $D$ i.e. sur $\partial\Delta_D$, sont les arcs de spécialisation \uline{et leurs antipodiques}, organisés en polygône de la façon habituelle. On considère que par l'un ou par l'autre, on peut aller vers une position rel. \unsure{adjacentes|adjacente}.

% @ R9–15
La surface $\unsure{\mathscr X|X|\mathcal X}$ obtenue, immergée dans $X$ par une immersion, \deleted{\unsure{il|et}} l'image inverse \inserted{$L$} de $X_1$ \deleted{est} donne lieu à \deleted{une} l'interprétation habituelle \deleted{de} \unsure{la|fé.} \unsure{p.v.|pos.|p.r.} (stable) associée à une face de $(\mathscr X,K)$.

% @ R16–end, heading and list with brace
Description combinatoire

\[
\begin{aligned}
\text{1) graphe }\Gamma\quad\left\{\begin{aligned}
&\text{sommets : pos. stables }D\\
&\text{arêtes orientées : paires }(D,a),\ a\text{ un}\\
&\deleted{\text{ori.}}\qquad\text{arc de transition sur }D\\
&\qquad\text{(de spécialisation,}\\
&\qquad\text{ou antipodique}\\
&\qquad\text{d'un tel)}
\end{aligned}\right.
\end{aligned}
\]
